\documentclass[10pt,a4paper]{amsart}
\usepackage[margin=19mm]{geometry}
\usepackage{amsmath,amssymb,mathtools}
\usepackage[scr=rsfs,cal=euler]{mathalfa}
\usepackage{xfrac}
\usepackage[unicode,hidelinks,bookmarks=false]{hyperref}
\newcommand{\ie}{i.e.\ }
\newcommand{\cf}{cf.\ }
\newcommand{\resp}{respectively\ }
\newcommand{\Subsection}{\subsection}
\providecommand{\nobreakdash}{\nobreak\hbox{-}}
\setlength{\emergencystretch}{2em}
\multlinegap0pt
\usepackage{amssymb}
\usepackage{mathtools}
\usepackage{tikz}
\usepackage{xcolor}
\newtheorem{prop}{Proposition}
\title{\textbf{Iterating the Lehmer code on inversion sequences: Catalan fixed points and finite stabilization}}
\author{Julian Allagan, Shanzhen Gao, Benjamin Testart}
\date{Source: arXiv:2608.24476v2 (CC BY 4.0)}
\begin{document}
\maketitle

\noindent\emph{Source and licence.} \textbf{Iterating the Lehmer code on inversion sequences: Catalan fixed points and finite stabilization}. Version arXiv:2608.24476v2 (CC BY 4.0), distributed by arXiv under the Creative Commons Attribution 4.0 International licence, http://creativecommons.org/licenses/by/4.0/, as recorded in the arXiv licence metadata for this exact version and shown on the abstract page https://arxiv.org/abs/2608.24476v2. Full licence notice: \texttt{licenses/arxiv-2608.24476-CC-BY-4.0.txt}.\par\smallskip

\noindent\textit{Editorial scope.} Counted unit: the article's prop prop:preimage\_characterization (source0.tex lines 273--277). The article's complete original proof of it is delivered with it (source0.tex lines 279--305, one \texttt{proof} environment of 27 lines). No mathematics is written, repaired or re-derived: the statement and the proof are the article's own lines, in the article's own notation, and no retained line is substituted or re-flowed.  Retained uncounted as the source-local context the proof needs: lines 142--144; lines 263--265.  Nothing was omitted from any retained span, so the package carries no omission record.  No retained line contains a citation, so the wrapper carries no bibliography and no bibliography entry.  No retained line contains a figure environment, an image-inclusion command or drawing code, so no image is delivered and the package declares no asset dependency.  Editorial formatting only, in addition: an AMS \texttt{amsart} preamble that restates the article's own declarations and macros used by the retained lines, namely the environments \texttt{prop} on separate counters, together with the standard packages those lines need.  The retained environments print in this document as Proposition 1 in order of appearance, whereas the article prints the same items with the numbering of its own full document.  The complete native source of the article as distributed in the arXiv e-print for this exact version is bundled unchanged as \texttt{sources/208481/source0.tex}.
\begin{prop}\label{prop:theta_basic}
For every finite integer sequence $\sigma$, one has $\Theta(\sigma)\in I_{|\sigma|}$. Moreover, for every $n\geq0$, the restriction $\Theta:S_n\to I_n$ is a bijection.
\end{prop}

\begin{prop}\label{prop:image_leftmax_saturated}
Let $\sigma$ be a nonempty integer sequence, let $p$ be the position of the leftmost maximum of $\sigma$, and let $\tau=\Theta(\sigma)$. Then $p$ is the position of the rightmost saturated entry of $\tau$.
\end{prop}

\begin{prop}\label{prop:preimage_characterization}
Let $\sigma$ and $\sigma'$ be two integer sequences of length $n \geq 0$. Then $\Theta(\sigma)=\Theta(\sigma')$ if and only if,
for every pair $p<q$, one has $\sigma_p<\sigma_q$ if and only if
$\sigma'_p<\sigma'_q$.
\end{prop}

\begin{proof}
Suppose first that $\sigma$ and $\sigma'$ have the same increasing pairs. Then, for each position $q$, the subsets $\{p<q:\sigma_p<\sigma_q\}$ and $\{p<q:\sigma'_p<\sigma'_q\}$ of $[1,q-1]$ are equal. Hence they have the same cardinality, and therefore $\Theta(\sigma)=\Theta(\sigma')$.

Conversely, assume $\Theta(\sigma)=\Theta(\sigma')=\tau$. We prove by induction on $n$ that $\sigma$ and $\sigma'$ have the same increasing pairs. The assertion is immediate for $n\leq1$, so let $n\geq2$.

Let $p$ be the position of the leftmost maximum of $\sigma$. By
Proposition~\ref{prop:image_leftmax_saturated}, $p$ is the rightmost saturated
position of $\tau$. Applying the same proposition to $\sigma'$ shows that $p$ is
also the position of the leftmost maximum of $\sigma'$. We first compare all pairs involving position $p$. If $i<p$, then $\sigma_i<\sigma_p$ and $\sigma'_i<\sigma'_p$, since $p$ is the leftmost maximum in both sequences. If $j>p$, then $\sigma_j\leq\sigma_p$ and $\sigma'_j\leq\sigma'_p$, since $p$ is a maximum in both sequences. Thus, for every $r\neq p$, the comparison between positions $p$ and $r$ has the same truth value in $\sigma$ as in $\sigma'$. Delete position $p$ from $\sigma,\sigma'$, and $\tau$, and denote the resulting sequences by $\widehat{\sigma},\widehat{\sigma}'$, and $\widehat{\tau}$, respectively. We claim that $\Theta(\widehat{\sigma})=\widehat{\tau}=\Theta(\widehat{\sigma}')$.

We verify the $\Theta$-relations after deletion. For a position originally
to the left of $p$, neither the entry nor any earlier entry is affected by the
deletion, so the corresponding $\Theta$-coordinate is unchanged. For a position
$q>p$, the deleted entry $\sigma_p$ satisfies $\sigma_p\geq\sigma_q$ and hence
does not contribute to $\Theta(\sigma)_q$; removing it therefore deletes no index
from the set counted by $\Theta(\sigma)_q$. After reindexing, the corresponding
coordinate of $\Theta(\widehat{\sigma})$ equals $\tau_q$. Hence
$\Theta(\widehat{\sigma})=\widehat{\tau}$, and in particular
$\widehat{\tau}\in I_{n-1}$ by Proposition~\ref{prop:theta_basic}. The identical
argument applied to $\sigma'$ gives $\Theta(\widehat{\sigma}')=\widehat{\tau}$.

By the induction hypothesis, $\widehat{\sigma}$ and $\widehat{\sigma}'$ have the
same increasing pairs. Together with the comparisons involving position $p$
verified above, this shows that $\sigma_a<\sigma_b$ if and only if
$\sigma'_a<\sigma'_b$ for every pair $a<b$ of original positions. Hence $\sigma$
and $\sigma'$ have the same increasing pairs, completing the proof.
\end{proof}

\end{document}
